\section{A quantitative exterior seminorm}
\label{sec:seminorm}

We prove the estimate in Theorem~\ref{thm:intro-estimate}.
Throughout this section $M$ is closed, oriented and negatively
curved, of dimension $n\geq2$. Write $B=M\times M$,
$G=\Gamma\times\Gamma$, $N=2n$, and
\[
 V_k=(\bigwedge^k I_{\R})_G,\qquad I_{\R}=I\otimes\R.
\]

We first compare the exterior seminorm with a cellular quotient seminorm.

For an ordered tuple of labels put
\[
 u(a_0,\ldots,a_k)
 =[(a_1-a_0)\wedge\cdots\wedge(a_k-a_0)]
 \in V_k.
\]
Define
\begin{equation}\label{eq:delta-norm}
 \|z\|_\Delta=
 \inf\left\{\sum_{j=1}^{r}|t_j|:
 z=\sum_{j=1}^{r}t_j
    u(a_{j,0},\ldots,a_{j,k})\right\}.
\end{equation}
The infimum is over finite expressions. This is a
possibly degenerate seminorm, not an asserted norm.

Let $Q=|Y|/G$, where $Y$ is the full simplex on
$\Gamma$. Endow the cellular chain complex of
$Q$ with the quotient $\ell^1$-seminorm assigning
cost one to each oriented simplex orbit.

\begin{lemma}\label{lem:isometric-image}
The injection
$\iota_k:H_k(Q;\R)\hookrightarrow V_k$
from Lemma~\ref{lem:cycle-injective} preserves
the quotient $\ell^1$-seminorm:
\[
 \|\iota_k[z]\|_\Delta
 =\inf\{\|c\|_1:\partial c=0,\ [c]=[z]\}.
\]
In particular, if $\beta:\Gamma^{k+1}\to\R$
is a bounded alternating $G$-invariant cocycle,
it induces a linear functional on $V_k$, and
\begin{equation}\label{eq:dual-pairing}
 |\langle\beta,z\rangle|
 \leq\|\beta\|_\infty\|z\|_\Delta.
\end{equation}
\end{lemma}

\begin{proof}
By Lemma~\ref{lem:exterior-quotient},
$V_k$ is the quotient of $C_k(Y;\R)_G$
by $\partial C_{k+1}(Y;\R)_G$. The generators
in \eqref{eq:delta-norm} are exactly the
images of the oriented simplex generators.
If $z$ is a cycle, every chain with the same
image in this quotient differs from it by a
boundary, and therefore is again a cycle.
Thus the infima defining both seminorms agree.

An alternating invariant cocycle vanishes
on cellular boundaries, so it defines a
functional on the quotient, of norm bounded
by its supremum norm. This proves
\eqref{eq:dual-pairing}.
\end{proof}

\begin{lemma}\label{lem:upper-norm}
For the class
$A_M=\mathfrak a_N\frown[B]$ one has
\[
 \|A_M\|_\Delta\leq\|B\|.
\]
\end{lemma}

\begin{proof}
Use a singular fundamental cycle of $B$ whose
$\ell^1$-norm is arbitrarily close to $\|B\|$.
The bar comparison followed by $T_N$ in
Proposition~\ref{prop:comparison} sends each
oriented singular simplex to an exterior
generator with coefficient of absolute
value at most one. Therefore the induced
chain map is nonexpanding in the $\ell^1$-norm.
Taking the infimum proves the inequality.
\end{proof}

The lower estimate is obtained from a bounded orientation cocycle.

By the bounded-cohomological dual description of
simplicial volume \cite{Gromov1982}, for every
$\epsilon>0$ there is a bounded alternating
homogeneous $n$-cocycle $\omega$ on $\Gamma$
representing the real orientation class,
normalized by
\begin{equation}\label{eq:orientation-norm}
 \langle[\omega],[M]\rangle=1,\qquad
 \|\omega\|_\infty\leq\|M\|^{-1}+\epsilon.
\end{equation}
Negative curvature ensures $\|M\|>0$.
Pull $\omega$ back to $G$ along the two
projections and denote the resulting cocycles by
$A=p_1^*\omega$ and $B'=p_2^*\omega$.
We write $B'$ to distinguish this cocycle
from the manifold $B=M\times M$.

Assume first that $n$ is odd.
Consider the $N$-cocycle on $G$
\[
 \alpha=\frac12(A\smile B'-B'\smile A).
\]
Its restriction to $\Delta\Gamma$ is identically
zero as a cochain. The cup-product sup-norm
estimate gives
\[
 \|\alpha\|_\infty\leq\|\omega\|_\infty^2.
\]
Since $n$ is odd, graded commutativity
in ordinary cohomology gives
\[
 \langle[\alpha],[M\times M]\rangle=1.
\]
By Theorem~\ref{thm:bounded-lift}, $\alpha$
lifts to a bounded alternating invariant
$N$-cocycle $\beta$ on $\Gamma^{N+1}$, with
\[
 \|\beta\|_\infty
 \leq N(N+1)!\,\|\omega\|_\infty^2.
\]
The bar comparison of Proposition
\ref{prop:comparison}, followed by
\eqref{eq:dual-pairing}, gives
\[
 1=|\langle\beta,A_M\rangle|
 \leq N(N+1)!\|\omega\|_\infty^2
       \|A_M\|_\Delta.
\]
Letting $\epsilon\downarrow0$ yields
\begin{equation}\label{eq:odd-bound}
 \|A_M\|_\Delta\geq
 \frac{\|M\|^2}{2n(2n+1)!}.
\end{equation}

We turn to the case where $n$ is even and set
\[
 D=B'-A\in C_b^n(G).
\]
As cochains, $D|_{\Delta\Gamma}=0$.
Theorem~\ref{thm:bounded-lift} supplies a
bounded alternating invariant cocycle
$\widetilde D$ on $\Gamma^{n+1}$
with $\pi^*[\widetilde D]=[D]$.

For the explicit bound we use the retraction
$r:G\to\Delta\Gamma$, $(g,h)\mapsto(g,g)$,
in the initial prism of Lemma
\ref{lem:relative-primitive}. In that prism,
the summands involving $p_1^*\omega$ vanish:
at the transition between $r(g_i,h_i)$ and
$(g_i,h_i)$ the first-coordinate arguments
repeat, and $\omega$ is alternating.
Hence only the $p_2^*\omega$ terms contribute.
The primitive therefore has norm at most
$n\|\omega\|_\infty$, rather than
$2n\|\omega\|_\infty$.
The remaining horizontal and vertical
contractions give
\begin{equation}\label{eq:Dtilde-bound}
 \|\widetilde D\|_\infty
 \leq n(n+1)!\|\omega\|_\infty.
\end{equation}

Define a bounded alternating invariant
$N$-cocycle on $\Gamma$ by
\[
 \beta=-\frac12
 \operatorname{Alt}(\widetilde D\smile\widetilde D),
\]
where $\operatorname{Alt}$ denotes normalized
alternation (average over permutations with
sign). Since normalized alternation has norm
at most one and induces the identity in
cohomology with real coefficients,
\[
 \|\beta\|_\infty
 \leq\frac12 n^2((n+1)!)^2
       \|\omega\|_\infty^2.
\]
The cohomology class of $\pi^*\beta$ equals
$-\frac12[D]^2$. In ordinary cohomology
of $M\times M$, dimensional reasons give
$[A]^2=[B']^2=0$ and, since $n$ is even,
$[A]\smile[B']=[B']\smile[A]$. Consequently
\[
 -\frac12[D]^2=[A]\smile[B']
 \quad\text{in }H^{2n}(M\times M;\R),
\]
so $\langle\pi^*[\beta],[M\times M]\rangle=1$.
The same bar comparison and duality as before
now yield
\begin{equation}\label{eq:even-bound}
 \|A_M\|_\Delta\geq
 \frac{2\|M\|^2}{n^2((n+1)!)^2}.
\end{equation}

\begin{proof}[Proof of Theorem~\ref{thm:intro-estimate}]
The upper bound is Lemma~\ref{lem:upper-norm}.
The lower bound is \eqref{eq:odd-bound}
when $n$ is odd and \eqref{eq:even-bound}
when $n$ is even. Since $\|M\|>0$,
the class $A_M$ is nonzero. Thus
$\mathfrak a_{2n}$ is nonzero as well.
\end{proof}

\begin{remark}
The constants are explicit rather than optimal.
Our argument does not show
$\|A_M\|_\Delta=\|M\times M\|$.
In particular, it does not prove equality
between the cellular seminorm of the nerve
and the simplicial volume of the product.
\end{remark}
